% Abhakseite „Das kann ich“ – Zone nur im Gesamt (\mitzone)
%% TODO Ich-kann-Sätze: Platzhalter wie in den Aufgabendateien
\begin{abhakseite}
\ifdefined\mitzone
\abhakgruppe{Kennst du schon}
\abhak{1}{Relative Häufigkeit und Anteil aus einer Stichprobe bilden}
\abhak{2}{Binomialverteilung ansetzen und Punktwahrscheinlichkeiten berechnen}
\abhak{3}{Graphen zweier Funktionen ablesen (Schnitt mit Waagerechten und Senkrechten)}
\abhak{4}{Die Sigma-Regeln und ihre c-Werte}
\abhak{5}{Quadratische Gleichungen und Wurzelgleichungen lösen, Lösungen deuten}
\abhak{6}{Ungleichungen in n lösen und runden}
\abhak{7}{Binomialverteilung ansetzen und Punktwahrscheinlichkeiten berechnen – Fehler finden}
\abhak{8}{Binomialverteilung ansetzen und Punktwahrscheinlichkeiten berechnen – selbst rechnen}
\fi
\abhakgruppe{Einheit 1 · Das Intervall lesen und deuten: die Überdeckungsdeutung (verträglich heißt: der angenommene Anteil wird vom Intervall überdeckt}
\abhak{9}{Lesen und Überdeckung}
\abhak{10}{Lesen und Überdeckung – Fehler finden, Begründen}
\abhakgruppe{Einheit 2 · Mit der Näherungsformel rechnen: die Grenzgleichung nach p lösen (quadratisch}
\abhak{11}{Näherungsformel}
\abhak{12}{Verträglichkeit einer vermuteten Trefferwahrscheinlichkeit über das 1,96σ-Intervall um den Erwartungswert prüfen}
\abhak{13}{Näherungsformel – Fehler finden, Begründen}
\abhakgruppe{Einheit 3 · Argumente über Formel und Verträglichkeit: die Länge des Intervalls bei doppeltem Umfang (Faktor eins durch Wurzel zwei}
\abhak{14}{Argumente}
\abhak{15}{Argumente – Fehler finden, Begründen}
\end{abhakseite}
